\documentclass{ctexart}

% 用于导入图片
% https://tex.stackexchange.com/a/23076
\usepackage{graphicx}
% 用于页面设置
\usepackage[a4paper]{geometry}
% 用于获取标题，作者等信息
\usepackage{titling}
% 用于下划线
\usepackage{ulem}
% 用于计算文字宽度
\usepackage{calc}
% 用于获得灰度模型
\usepackage{xcolor}
% 用于数学公式
\usepackage{amsmath}
\usepackage{lstfiracode} % https://ctan.org/pkg/lstfiracode
% 设置代码块风格
\lstset{
  language=Matlab, 
  numbers=left,
  backgroundcolor=\color{lightgray!40!white},
  basicstyle=\ttfamily,   % Use \ttfamily for source code listings
  keywordstyle=\color{blue!70}\bfseries,
  commentstyle=\ttfamily\color{green!40!black},
}
% 设置页面大小
\geometry{hmargin=3.17cm, vmargin=2.54cm}
\title{《数值分析》\\ 实\hspace{0.5cm}验\hspace{0.5cm}报\hspace{0.5cm}告}
\author{\pyvar{author_info.author_name.value}}
% 行间距
\linespread{1.3}

\ctexset{
    section={
        name={实验},
        format=\centering\kaishu\zihao{-2},
        number=\chinese{section},
    },
    subsection={
        name={,、},
        format=\kaishu\zihao{-4},
        number=\chinese{subsection},
    },
    subsubsection={
        name={,、},
        format=\kaishu\zihao{5},
        number=\arabic{subsubsection}
    }
}

% 重新定义标题页
\renewcommand{\maketitle}{
    \begin{titlepage}
        \begin{center}
            \includegraphics[width=2.84cm]{\pyvar{"imgs/schoolLogo1.png" | real_relative_path}}
        \end{center}
        \begin{center}
            \includegraphics[width=6.67cm]{\pyvar{"imgs/schoolLogo2.png" | real_relative_path}}
        \end{center}
        \zihao{-0}
        \heiti
        % 标题
        \begin{center}
            \thetitle
        \end{center}
        \vspace{1cm}

        \kaishu
        \zihao{2}

        \newlength{\underlineLength}
        \setlength{\underlineLength}{8cm}
        \newlength{\lrSpaceLength}

        %% for item in  ["author_name", "author_student_number", "author_major_name", "author_class_name"]
        \setlength{\lrSpaceLength}{
            (\underlineLength - \widthof{\pyvar{author_info[item].value}}) / 2
        }
        % \pyvar{author_info[item]["trans"]}
        \begin{center}
            \pyvar{author_info[item].trans}：\uline{\hspace{\lrSpaceLength}\pyvar{author_info[item].value}\hspace{\lrSpaceLength}}
        \end{center}
        %% endfor
        \vspace{1cm}
        \begin{center}
            \selectcolormodel{gray}
            \includegraphics[width=2.33cm]{\pyvar{"imgs/majorLogoGray.png" | real_relative_path}}
        \end{center}
        \begin{center}
            \zihao{-1}
            计算机科学学院 \\
            \vspace{1cm}
            \zihao{-2}
            2019-2020学年 第2学期
        \end{center}
    \end{titlepage}
}


\begin{document}
% 标题页
\maketitle

\section{}

\zihao{-5}

\subsection{问题的描述}


\subsubsection{设线性方程组}



\[
    \begin{bmatrix}
        4  & 2  & -3  & -1 & 2  & 1   & 0  & 0  & 0  & 0  \\
        8  & 6  & -5  & -3 & 6  & 5   & 0  & 1  & 0  & 0  \\
        4  & 2  & -2  & -1 & 3  & 2   & -1 & 0  & 3  & 1  \\
        0  & -2 & 1   & 5  & -1 & 3   & -1 & 1  & 9  & 4  \\
        -4 & 2  & 6   & -1 & 6  & 7   & -3 & 3  & 2  & 3  \\
        8  & 6  & -8  & 5  & 7  & 17  & 2  & 6  & -3 & 5  \\
        0  & 2  & -1  & 3  & -4 & 2   & 5  & 3  & 0  & 1  \\
        16 & 10 & -11 & -9 & 17 & 34  & 2  & -1 & 2  & 2  \\
        4  & 6  & 2   & -7 & 13 & 9   & 2  & 0  & 12 & 4  \\
        0  & 0  & -1  & 8  & -3 & -24 & -8 & 6  & 3  & -1
    \end{bmatrix}
    \begin{bmatrix}
        x_1 \\
        x_2 \\
        x_3 \\
        x_4 \\
        x_5 \\
        x_6 \\
        x_7 \\
        x_8 \\
        x_9 \\
        x_{10}
    \end{bmatrix}
    =
    \begin{bmatrix}
        5  \\
        12 \\
        3  \\
        2  \\
        3  \\
        46 \\
        13 \\
        38 \\
        19 \\
        -21
    \end{bmatrix}
\]

\[x^* = (1, -1, 0, 1, 2, 0, 3, 1, -1, 2) ^ T\]

利用 Gauss 顺序消去法或 Gauss 列主元消去法求解。

\subsubsection{设对称正定阵系数阵线方程组}



\[
    \begin{bmatrix}
        4  & 2  & -4 & 0  & 2   & 4  & 0   & 0  \\
        2  & 2  & -1 & -2 & 1   & 3  & 2   & 0  \\
        -4 & -1 & 14 & 1  & -8  & -3 & 5   & 6  \\
        0  & -2 & 1  & 6  & -1  & -4 & -3  & 3  \\
        2  & 1  & -8 & -1 & 22  & 4  & -10 & -3 \\
        4  & 3  & -3 & -4 & 4   & 11 & 1   & -4 \\
        0  & 2  & 5  & -3 & -10 & 1  & 14  & 2  \\
        0  & 0  & 6  & 3  & -3  & -4 & 2   & 19
    \end{bmatrix}
    \begin{bmatrix}
        x_1 \\
        x_2 \\
        x_3 \\
        x_4 \\
        x_5 \\
        x_6 \\
        x_7 \\
        x_8
    \end{bmatrix}
    =
    \begin{bmatrix}
        0   \\
        -6  \\
        20  \\
        23  \\
        9   \\
        -22 \\
        -15 \\
        45
    \end{bmatrix}
\]

\[
    x^* = (1, -1, 0, 2, 1, -1, 0, 2)^T
\]

利用平方根法或改进平方根法求解。

\subsubsection{三对角形线性方程组}



\[
    \begin{bmatrix}
        4  & -1 & 0  & 0  & 0  & 0  & 0  & 0  & 0  & 0  \\
        -1 & 4  & -1 & 0  & 0  & 0  & 0  & 0  & 0  & 0  \\
        0  & -1 & 4  & -1 & 0  & 0  & 0  & 0  & 0  & 0  \\
        0  & 0  & -1 & 4  & -1 & 0  & 0  & 0  & 0  & 0  \\
        0  & 0  & 0  & -1 & 4  & -1 & 0  & 0  & 0  & 0  \\
        0  & 0  & 0  & 0  & -1 & 4  & -1 & 0  & 0  & 0  \\
        0  & 0  & 0  & 0  & 0  & -1 & 4  & -1 & 0  & 0  \\
        0  & 0  & 0  & 0  & 0  & 0  & -1 & 4  & -1 & 0  \\
        0  & 0  & 0  & 0  & 0  & 0  & 0  & -1 & 4  & -1 \\
        0  & 0  & 0  & 0  & 0  & 0  & 0  & 0  & -1 & 4
    \end{bmatrix}
    \begin{bmatrix}
        x_1 \\
        x_2 \\
        x_3 \\
        x_4 \\
        x_5 \\
        x_6 \\
        x_7 \\
        x_8 \\
        x_9 \\
        x_{10}
    \end{bmatrix}
    =
    \begin{bmatrix}
        7   \\
        5   \\
        -13 \\
        2   \\
        6   \\
        -12 \\
        -14 \\
        -4  \\
        5   \\
        -5
    \end{bmatrix}
\]

利用追赶法求解。}

\subsection{方法描述}

\subsubsection{高斯消去法}

高斯消去法，该方法以数学家卡尔·高斯命名，但最早出现于中国古籍《九章算术》，成书于约公元前150年.

设有线性方程组：

\begin{equation}
    \left\{\begin{aligned}
         & a_{11}x_1 +  a_{12}x_2  +  \cdots  +  a_{1n}x_n = b_1, \\
         & a_{21}x_1 +  a_{22}x_2  +  \cdots  +  a_{2n}x_n = b_2, \\
         & \qquad\qquad\qquad\qquad\qquad\qquad\quad\vdots        \\
         & a_{m1}x_1 +  a_{m2}x_2  +  \cdots  +  a_{mn}x_n = b_m.
    \end{aligned}\right.
    \label{线性方程组示例}
\end{equation}

或写成矩阵形式：

\begin{equation*}
    \begin{bmatrix}
        a_{11} & a_{12} & \cdots & a_{1n} \\
        a_{21} & a_{22} & \cdots & a_{2n} \\
        \vdots & \vdots &        & \vdots \\
        a_{m1} & a_{m2} & \cdots & a_{mn}
    \end{bmatrix}
    \begin{bmatrix}
        x_1    \\
        x_2    \\
        \vdots \\
        x_n
    \end{bmatrix}
    =
    \begin{bmatrix}
        b_1    \\
        b_2    \\
        \vdots \\
        b_m
    \end{bmatrix}
\end{equation*}

简记为 $Ax = b$.

用高斯消去法解方程组的基本思想是用逐次消去未知数的方法把源方程组 $Ax = b$ 化为与其等价的三角形方程组，而求解三角形方程组可用回代的方法求解. 换句话来说，用行的初等变换将原方程组的系数矩阵化为简单形式（上三角矩阵），从而将求解原方程组 \eqref{线性方程组示例} 的问题转化为求解简单方程组的问题，或者说，对系数矩阵 $A$ 实行一些左变换（用一些简单矩阵）将其约化为上三角矩阵.

设 $Ax = b$，其中 $A \in R^{m \times n}$.

\begin{enumerate}
    \item 如果 $a_{kk}^{k} \neq 0\quad(k=1, 2 , \cdots, n)$，则可以通过高斯消去法将 $Ax = b$ 约化为等价的三角形方程组，且计算公式为：
          \begin{enumerate}
              \item 消元计算$(k=1, 2, \cdots, n - 1)$
                    \begin{equation*}
                        \left\{\begin{aligned}
                             & m_{ik}         = a_{ik}^{(k)} / a_{kk}^{(k)} \quad (i = k + 1, \cdots, n),          \\
                             & a_{ij}^{(k+1)} = a_{ij}^{(k)} - m_{ik}a_{kj}^{(k)} \quad (i, j = k + 1, \cdots, n), \\
                             & b_i^{(k+1)}    = b_i^{(k)} - m_{ik}b_k^{(k)} \quad (i = k + 1, \cdots, n).
                        \end{aligned}\right.
                    \end{equation*}
              \item 回代计算
                    \begin{equation*}
                        \left\{\begin{aligned}
                             & m_n = b_n^{(n)} / a_{nn}^{(n)},                                                                         \\
                             & x_i = (b_i^{(i)} - \sum_{j = i + 1}^n a_{ij}^{(i)}x_j) / a_{ii} ^{(i)} \quad (i = n - 1, \cdots, 2, 1).
                        \end{aligned}\right.
                    \end{equation*}
          \end{enumerate}
    \item 如果 $A$ 为非奇异矩阵，则可通过高斯消去法（及交换两行的初等变换）将方程组 $Ax = b$ 进行约化
\end{enumerate}

\subsubsection{平方根法}

平方根法，也称为 choleskly 方法，用于将给定的系数矩阵 \(A\) 分解为
\(A = L L^T\) 的形式.这里的系数矩阵 \(A\) 是对称正定矩阵，其中的 \(L\)
是下三角矩阵.

由此可以写出以下等式：

\begin{equation*}
    \begin{aligned}
        A = LL^T & =
        \begin{bmatrix}
            l_{11} & 0      & 0      \\
            l_{21} & l_{22} & 0      \\
            l_{31} & l_{32} & l_{33} \\
        \end{bmatrix}
        \begin{bmatrix}
            l_{11} & l_{21} & l_{31} \\
            0      & l_{22} & l_{32} \\
            0      & 0      & l_{33} \\
        \end{bmatrix} \\
                 & =
        \begin{bmatrix}
            l_{11}^2     &                             & (symmetric)                     \\
            l_{21}l_{11} & l_{21}^2 + l_{22}^2                                           \\
            l_{31}l_{11} & l_{31}l_{21} + l_{32}l_{22} & l_{31}^2 + l_{32}^2 + l_{33} ^2 \\
        \end{bmatrix}
    \end{aligned}
\end{equation*}

由此可以得到：

\begin{equation*}
    L =
    \begin{bmatrix}
        (\pm) \sqrt{a_{11}} & 0                            & 0                                     \\
        a_{21} / l_{11}     & (\pm)\sqrt{a_{22}-l_{21}^2}  & 0                                     \\
        a_{31} / l_{11}     & (a_{32}-l_{31}l_{21})/l_{22} & (\pm) \sqrt{a_{33}-l_{31}^2-l_{32}^2} \\
    \end{bmatrix}
\end{equation*}

根据规律可以归纳得到：

\begin{equation*}
    l_{j,j} = (\pm) \sqrt{a_{j,j} - \sum_{k=1}^{j-1} l_{j,k}^2} \\
\end{equation*}

\[
    l_{i,j} = (a_{i,j} - \sum_{k=1}^{j-1}l_{i,k}l_{j,k}) / l_{j,j},  i > j
\]

因此，要计算${l}_{i,j}(i\neq j)$ 只需利用矩阵 $L$ 的左、上方元素的值.计算通常是以下面其中一种顺序进行的.\footnote{https://zh.wikipedia.org/wiki/Cholesky}

\begin{enumerate}
    \item Cholesky–Banachiewicz：算法从矩阵L的左上角开始，依行进行计算.
    \item Cholesky–Crout：算法从矩阵L的左上角开始，依列进行计算.
\end{enumerate}

若有需要，整个矩阵可以逐个元素计算得出，无论使用何种顺序读取.

将系数矩阵 \(A\) 分解为 \(A = L L^T\) 后，由 \(Ax = b\)
可得：\(LL^Tx = b\).

如果进行以下分解：

\[
    \left\{\begin{aligned}
         & y  = L^Tx \\
         & Ly = b
    \end{aligned}\right.
\]


就可以先计算 \(y\)，再计算 \(x\)，因为 \(L\) 是下三角矩阵，而 \(L^T\)
是上三角矩阵，这样会比直接使用 \(A\) 计算更简便.

\subsubsection{追赶法}

设系数矩阵 $A$ 为对角占优的三对角线矩阵，则有：

\begin{equation}
    A =
    \begin{bmatrix}
        b_1 & c_1                                  \\
        a_2 & b_2    & c_2                         \\
            & \ddots & \ddots  & \ddots            \\
            &        & a_{n-1} & b_{n-1} & c_{n-1} \\
            &        &         & a_n     & b_n
    \end{bmatrix}
    \begin{bmatrix}
        x_1     \\
        x_2     \\
        \vdots  \\
        x_{n-1} \\
        x_n
    \end{bmatrix}
    =
    \begin{bmatrix}
        f_1     \\
        f_2     \\
        \vdots  \\
        f_{n-1} \\
        f_n
    \end{bmatrix}
    \label{追赶法:三对角线方程组}
\end{equation}

简记为 \textbf{$Ax=f$}. 其中当 $\left|i-j\right| > 1$ 时，$a_{ij} = 0$，且有：

\begin{enumerate}
    \item $\left|b_1\right| > \left|c_1\right| > 0$；
    \item $\left|b_i\right| \geq \left|a_i\right| + \left|c_i\right|, a_i, c_i \neq 0\quad(i=2,3,\cdots,n-1)$；
    \item $\left|b_n\right| > \left|a_n\right| > 0$.
\end{enumerate}

利用矩阵的直接三角分解法来推导解三对角线方程组 \eqref{追赶法:三对角线方程组} 的计算公式，由于系数阵 $A$ 的特点，可以将 $A$ 分解为两个三角阵的乘积，即:

\[A=LU\]

其中 $L$ 为下三角矩阵，$U$ 为单位上三角矩阵，这种分解是可能的，设：

\begin{equation}
    A =
    \begin{bmatrix}
        b_1 & c_1                                  \\
        a_2 & b_2    & c_2                         \\
            & \ddots & \ddots  & \ddots            \\
            &        & a_{n-1} & b_{n-1} & c_{n-1} \\
            &        &         & a_n     & b_n
    \end{bmatrix}
    =
    \begin{bmatrix}
        \alpha_1                                  \\
        \gamma_2 & \alpha_2                       \\
                 & \ddots   & \ddots              \\
                 &          & \gamma_n & \alpha_n
    \end{bmatrix}
    \begin{bmatrix}
        1 & \beta_1                        \\
          & 1       & \ddots               \\
          &         & \ddots & \beta_{n-1} \\
          &         &        & 1
    \end{bmatrix}
    \label{追赶法:分解}
\end{equation}

其中 $\alpha_i$，$\beta_i$，$r_i$ 为待定系数，比较 \eqref{追赶法:分解} 两边可得

\begin{equation}
    \left.\begin{aligned}
        b_1      = \alpha_1, c_1 = \alpha_1\beta_1,                                       & \\
        \alpha_i = \gamma_i, b_i = \gamma_i\beta_{i-1} + \alpha_i\quad(i = 2, \cdots, n), & \\
        c_i      = \alpha_i\beta_i\quad(i=2,3,\cdots,n-1).                                &
    \end{aligned}\right\}
    \label{追赶法:推导公式}
\end{equation}

由 $\alpha_1 = b_1 \neq 0$，$\left|b_1\right| > \left|c_1\right| > 0$，$\beta_1 = c_1 / b_1$，得 $0 < \left|\beta_1\right| < 1$. 用归纳法可以证明

\[ \left|\alpha_i\right| > \left|c_i\right| \neq 0\quad(i=1,2, \cdots, n - 1)\]

即 $0 < \left|\beta_i\right| < 1$，从而可以由 \eqref{追赶法:推导公式} 求出 $\beta_i$.

求解 $Ax=f$ 等价于解两个三角形方程组

\begin{enumerate}
    \item $Ly=f$，求 $y$；
    \item $Ux=y$，求 $x$.
\end{enumerate}

从而得到解三对角线方程组的追赶法公式：

\begin{enumerate}
    \item 计算 $\{\beta_i\}$ 的递推公式 \\
          $\beta_1 = c_1 / b_1$，
          $\beta_i = c_i / (b_i - a_i\beta_{i-1})$，$\quad(i=2,3,\cdots,n-1)$；
    \item 解 $Ly = f$ \\
          $y_1 = f_1 / b_1$，
          $y_i = (f_i - a_iy_{i-1}) / (b_i - a_i\beta_{i-1})\quad(i=2,3,\cdots,n)$；
    \item 解 $Ux = y$ \\
          $x_n = y_n$，
          $x_i = y_i - \beta_ix_{i+1}\quad(i=n-1,n-2,\cdots,2,1)$.
\end{enumerate}

计算系数 $\beta_1 \rightarrow \beta_2 \rightarrow \cdots \rightarrow \beta_{n-1}$ 以及 $y_1 \rightarrow y_2 \rightarrow \cdots \rightarrow y_n$的过程称为追的过程，将计算方程组的解 $x_n \rightarrow x_{n-1} \rightarrow \cdots \rightarrow x_1$ 的过程称为赶的过程.
}

\subsection{算法设计}

\subsubsection{问题一}

\begin{breakablealgorithm}
    \caption{高斯顺序消去法}
    \begin{algorithmic}[1]
        \REQUIRE{矩阵 $A$ 的对角线元素均不为 0}
        \STATE $B \leftarrow [A, b]$ \COMMENT{通过合并系数矩阵 $A$ 和常数矩阵 $b$ 获得增广矩阵 $B$ }
        \STATE $[n, m] \leftarrow size(B)$ \COMMENT{求增广矩阵 B 的行数和列数，分别记做 n 和 m}
        \FOR{$k=1$ \TO $n$}
        \FOR{$o=k+1$ \TO $n$}
        \FOR{$p=0$ \TO $m$}
        \STATE $B_{o, p} = B_{o, p} - (B_{o, k} / B_{k, k}) * B_{k, p}$
        \ENDFOR
        \ENDFOR
        \ENDFOR
        \STATE $x \leftarrow zeros(m-1, 1)$ \COMMENT{初始化 x 向量}
        \FOR{$k=n$ \TO $1$ by step $-1$}
        \FOR{$q=k+1$ \TO $m-1$}
        \STATE $x_k \leftarrow (B_{k, m} - B_{k, q} * x_q ) / B_{k, k}$
        \ENDFOR
        \ENDFOR
    \end{algorithmic}
\end{breakablealgorithm}

\subsubsection{问题二}

\begin{breakablealgorithm}
    \caption{cholesky 算法}
    \begin{algorithmic}[1]
        \label{alg:cholesky}
        \REQUIRE{被分解的矩阵必须为厄米特矩阵或符合正定矩阵定义的矩阵}
        \STATE $[n, \sim] \leftarrow size(A)$ \COMMENT{获取矩阵 A 的行数 n}
        \STATE $L \leftarrow zeros(n, n)$ \COMMENT{初始化矩阵 L}
        \FOR{$j=1$ \TO $n$}
        \FOR{$i=j$ \TO $n$}
        \IF{$i == j$}
        \STATE $L_{i, j} \leftarrow \sqrt{A_{i, j}} - \sum_{k=1}^{j-1}L_{j, k}^2$
        \ELSE
        \STATE $L_{i, j} \leftarrow (A_{i, j} - \sum_{k=1}^{j-1}L_{i, k}L_{j, k}) / L_{j, j}$
        \ENDIF
        \ENDFOR
        \ENDFOR
    \end{algorithmic}
\end{breakablealgorithm}

\begin{breakablealgorithm}
    \caption{回代算法}
    \begin{algorithmic}[1]
        \STATE $[n, \sim] \leftarrow size(A)$ \COMMENT{获取矩阵 A 的行数 n}
        \STATE $[L, U] \leftarrow cholesky(A)$ \COMMENT{通过算法 \ref{alg:cholesky} 分解系数矩阵 A}
        \STATE $y = zeros(n, 1)$ \COMMENT{初始化 y}
        \FOR{$k = 1$ \TO $n$}
        \STATE $y_k = (b_k - \sum_{j = 1}^{k - 1} L_{k, j}y_j) / L_{k, k}$
        \ENDFOR
        \STATE $x = zeros(n, 1)$ \COMMENT{初始化 x}
        \FOR{$k = n$ \TO 1 step by -1}
        \STATE $x_k = (y_k - \sum_{j=k + 1}^n U_{k, j} * x_j) / U_{k, k}$
        \ENDFOR
    \end{algorithmic}
\end{breakablealgorithm}

\subsubsection{问题三}

\begin{breakablealgorithm}
    \caption{追赶法}
    \begin{algorithmic}[1]
        % \REQUIRE{系数矩阵 A 必须为三对角线矩阵}
        \STATE $[n, \sim] \leftarrow size(A)$ \COMMENT{获取矩阵 A 的行数 n}
        \STATE $a \leftarrow ones(n - 1, 1)$  \COMMENT{初始化 a 向量}
        \STATE $b \leftarrow ones(n, 1)$ \COMMENT{初始化 b 向量}
        \STATE $c \leftarrow ones(n - 1, 1)$\COMMENT{初始化 c 向量}
        \COMMENT{构造 a，b，c 向量}
        \FOR{$i=1$ \TO $n - 1$}
        \STATE $b_{i, 1} \leftarrow A_{i, i}$
        \STATE $a_{i, 1} \leftarrow A_{i + 1, i}$
        \STATE $c_{i, 1} \leftarrow A_{i, i + 1}$
        \ENDFOR
        \STATE $b_{n, 1} \leftarrow A_{n, n}$
        \STATE $d \leftarrow zeros(n - 1, 1)$ \COMMENT{初始化 d 向量（$\beta$ 向量）}
        \STATE $d_1 \leftarrow c_1 / b_1$
        \COMMENT{构造 d 向量}
        \FOR{$i=2$ \TO $n - 1$}
        \STATE $d_i \leftarrow c_i / (b_i - a_{i - 1} * d_{i - 1})$
        \ENDFOR
        \STATE $y \leftarrow zeros(n, 1)$ \COMMENT{初始化 y 向量}
        \STATE $y_1 \leftarrow f_1 / b_1$
        \FOR{$i=2$ \TO $n$}
        \STATE $y_i \leftarrow (f_i - a_{i - 1} * y_{i - 1}) / ((b_i - a_{i - 1}) * d_{i - 1})$
        \STATE $x \leftarrow zeros(n, 1)$ \COMMENT{初始化 x 向量}
        \STATE $x_n \leftarrow y_n$
        \ENDFOR
        \FOR{$i=n-1$ \TO $1$ by step $-1$}
        \STATE $x_i = y_i - d_i * x_{i + 1}$
        \ENDFOR
    \end{algorithmic}
\end{breakablealgorithm}
}

\subsection{程序代码}

%% macro include_by_name(names)
%% for name in names
\noindent
\textit{\pyvar{name + ".m"}}
\begin{lstlisting}[frame=single]
\pyblock{set rel_path = name+".m"}
\pyvar{rel_path | include_code}
\end{lstlisting}
%% endfor
%% endmacro

\subsubsection{问题一}
\pyvar{include_by_name(["gauss"])}
\subsubsection{问题二}
\pyvar{include_by_name(["cholesky", "squareRoot"])}
\subsubsection{问题三}
}

\subsection{计算结果及其分析}

\subsubsection{问题一}

\subsubsection{问题二}

\begin{lstlisting}[frame=single]
>> squareRoot(A, b)

ans =

    121.1481
    -140.1127
    29.7515
    -60.1528
    10.9120
    -26.7963
    5.4259
    -2.0185

\end{lstlisting}

利用平方根法求得的解为：

\[x = (121.1481, -140.1127, 29.7515, -60.1528, 10.9120, -26.7963, 5.4259, -2.0185)^T\]

虽然系数矩阵的对角元素都大于零，原则上可以不必选择主元，但由于矩阵的数值问题较大，不选择主元的结果就是产生较大的误差，所以在求解过程中还是应该选择主元以消除此误差，提高精度。

\subsubsection{问题三}
}

\subsection{结论}

在此次课设中，我通过使用高斯顺序消元法，平方根法和追赶法分别求解了三个线性方程组，并分析了三种方法在求解线性方程组的特点。

解线性方程组，应根据它自身的特点，选择最适合的算法来进行求解，没有一种算法是最好的，只能说在某种情形下，某种算法比另外一种算法更好。
}

\end{document}
